\documentclass[10pt,a4paper]{amsart}
\usepackage[margin=19mm]{geometry}
\usepackage{amsmath,amssymb,mathtools}
\usepackage[scr=rsfs,cal=euler]{mathalfa}
\usepackage{xfrac}
\usepackage[unicode,hidelinks,bookmarks=false]{hyperref}
\newcommand{\ie}{i.e.\ }
\newcommand{\cf}{cf.\ }
\newcommand{\resp}{respectively\ }
\newcommand{\Subsection}{\subsection}
\providecommand{\nobreakdash}{\nobreak\hbox{-}}
\setlength{\emergencystretch}{2em}
\multlinegap0pt
\usepackage{xcolor}
\usepackage{tcolorbox}
\usepackage{amssymb}
\usepackage[shortlabels]{enumitem}
\usepackage{tikz}
\newtheorem{theorem}{Theorem}
\title{Correctness of Extended RSA Public Key Cryptosystem}
\author{Dar-jen Chang, Suranjan Gautam}
\date{Source: arXiv:2512.24531v1 (CC BY 4.0)}
\begin{document}
\maketitle

\noindent\emph{Source and licence.} Correctness of Extended RSA Public Key Cryptosystem. Version arXiv:2512.24531v1 (CC BY 4.0), distributed by arXiv under the Creative Commons Attribution 4.0 International licence, http://creativecommons.org/licenses/by/4.0/, as recorded in the arXiv licence metadata for this exact version and shown on the abstract page https://arxiv.org/abs/2512.24531v1. Full licence notice: \texttt{licenses/arxiv-2512.24531-CC-BY-4.0.txt}.\par\smallskip

\noindent\textit{Editorial scope.} Counted unit: the article's theorem theorem (source0.tex lines 776--778). The article's complete original proof of it is delivered with it (source0.tex lines 779--803, one \texttt{proof} environment of 25 lines). No mathematics is written, repaired or re-derived: the statement and the proof are the article's own lines, in the article's own notation, and no retained line is substituted or re-flowed.  No further source-local context is needed: every cross-reference of every retained line resolves inside the two delivered spans.  Nothing was omitted from any retained span, so the package carries no omission record.  No retained line contains a citation, so the wrapper carries no bibliography and no bibliography entry.  No retained line contains a figure environment, an image-inclusion command or drawing code, so no image is delivered and the package declares no asset dependency.  Editorial formatting only, in addition: an AMS \texttt{amsart} preamble that restates the article's own declarations and macros used by the retained lines, namely the environments \texttt{theorem} on separate counters, together with the standard packages those lines need.  The retained environments print in this document as Theorem 1 in order of appearance, whereas the article prints the same items with the numbering of its own full document.  The complete native source of the article as distributed in the arXiv e-print for this exact version is bundled unchanged as \texttt{sources/191905/source0.tex}.
\begin{theorem}{Theorem}{Euler's Theorem.}{}{Euler-Theorem}%
For any two positive integers \(m \text{ and } N\), if \(\gcd(m,N) = 1\), then \(m^{\phi(n)} \equiv 1 \pmod N\).%
\end{theorem}

\begin{proof}{Proof}{}{Euler-Theorem-3}
The proof of Euler's Theorem is equvalant to prove for any \(m \in \phi\text{-set}(N)\),%
\begin{equation*}
m^{\phi(n)} \equiv 1 \pmod N\text{.}
\end{equation*}
%
\par
Since \(\phi\text{-set}(N)\) forms a group under multiplication modulo \(N\), for any \(m \in \phi\text{-set}(N)\), the cyclic group generated by \(m\) is a subgroup of \(\phi\text{-set}(N)\). Therefore, the order of \(m\), denoted by \(order(m)\), is a divisor of \(\phi(N)\), which is the order of the group \(\phi\text{-set}(N)\).%
\par
Therefore, we have%
\par
%
\begin{equation*}
m^{order(m)} \equiv 1 \pmod N
\end{equation*}
and%
\begin{equation*}
\phi(N) = k \times order(m), \text{ where } k \text{ is a positive integer.}
\end{equation*}
Combining the above results, we get%
\begin{equation*}
m^{\phi(N)} = m^{k \times order(m)} = (m^{order(m)})^k \equiv 1^k \equiv 1 \pmod N\text{.}
\end{equation*}
%
\end{proof}

\end{document}
